\section{The vertex-supported obstruction and the cell-pair construction}\label{sec:no-go}
The scalar mesh-refinement theorem does not reconstruct a continuum kernel.
Its vertex-supported candidate contributes zero to the continuum coupling
functional on an atomless substrate. This obstruction concerns that particular
representation. The development also provides a cell-pair construction which
converges to a specified continuous kernel.

\subsection{The target structure constrains nothing}
A \texttt{ContinuousNeuralField} over a substrate $M$ carries a drift
$\omega:M\to\mathbb{R}$, a kernel $K:M\times M\to\mathbb{R}$, and a timescale
$\tau\in\mathbb{R}$, with no conditions relating them.

\begin{claim}
For any $\omega$, $K$ and $\tau$ whatsoever, a \texttt{ContinuousNeuralField}
over $M$ exists.
\end{claim}
\begin{proof}
The three data are exactly the structure's fields, so the anonymous constructor
supplies the required inhabitant.
\end{proof}
\formal{continuousneuralfield-free}
\scope A constructor taking a mesh to such a field would therefore typecheck
with a kernel unrelated to that mesh, and its proof would be a definitional
unfolding. The statement records that the edge would carry no information, not
that a limit fails to converge.

\subsection{The candidate kernel does carry the weights}
Write $\mathrm{emb}$ for a triangulation's vertex embedding into $M$ and $w$ for
its edge weights. The induced kernel places each weight at its embedded pair,
\[
  \mathrm{vertexKernel}(x,y)=\sum_{u}\sum_{v}
    \bigl[\,x=\mathrm{emb}\,u \ \wedge\ y=\mathrm{emb}\,v\,\bigr]\;w(u,v),
\]
and vanishes off the image of the vertex set.

\begin{claim}
The kernel is sited on the finite set $\mathrm{emb}[V]$: if $x\notin
\mathrm{emb}[V]$ then $\mathrm{vertexKernel}(x,y)=0$ for every $y$.
\end{claim}
\begin{proof}
Fix such an $x$ and any $y$. Every summand carries the condition
$x=\mathrm{emb}\,u$; were it met for some $u$, $x$ would lie in the image. So
each summand is zero, and both sums vanish.
\end{proof}
\formal{vertexkernel-sitedon}

\begin{claim}
If $\mathrm{emb}$ is injective, then
$\mathrm{vertexKernel}(\mathrm{emb}\,u,\mathrm{emb}\,v)=w(u,v)$.
\end{claim}
\begin{proof}
Injectivity makes $\mathrm{emb}\,u=\mathrm{emb}\,b$ fail for $b\neq u$, so the
outer sum reduces to its $u$th term, and the inner sum to its $v$th by the same
argument. That term is $w(u,v)$.
\end{proof}
\formal{vertexkernel-apply-embedding}
\scope This is what stops the next result from being a statement about a kernel
that was zero from the outset. At embedded pairs the kernel takes precisely the
value the triangulation assigns.

\subsection{The no-go}
The field results are stated through the coupling functional
\[
  \mathrm{fieldCorrelation}(\mu,\theta,K)
    =\int_M\!\Bigl(\int_M K(x,y)\,\cos\bigl(\theta(y)-\theta(x)\bigr)\,d\mu(y)\Bigr)d\mu(x).
\]

\begin{claim}
On a substrate whose measure has no atoms,
$\mathrm{fieldCorrelation}(\mu,\theta,\mathrm{vertexKernel})=0$ --- for every
weight assignment and every phase field.
\end{claim}
\begin{proof}
By the siting claim the inner integral vanishes for each $x$ outside the finite
set $\mathrm{emb}[V]$. A finite set is a finite union of singletons, and no
singleton carries mass, so that set is null. The outer integral of a function
supported in a null set is zero.
\end{proof}
\formal{vertexkernel-fieldcorrelation-eq-zero}
\scope The weights are unconstrained and the phase field is arbitrary, so no
choice of triangulation data evades the conclusion. The obstruction is that the
discrete data occupies a null set, not that the weights are small.

\subsection{What this locates}
The coarse-graining theorem converges a discrete coupling \emph{energy} --- one
real number per triangulation --- to an integral over the substrate, and the
edge-weight sum reaches the same scalar. Both integrate the pair structure away,
while a continuum kernel is a function of two points. Spreading the weights over
cells rather than vertices would require a product-structured decomposition of
$M\times M$; the triangulation's edge regions are subsets of $M$ alone.

\subsection{The positive construction and its input}
\texttt{KernelMesh} supplies a finite measurable partition and one sample per
cell. It extends each sampled pair value over the whole cell pair.
\texttt{KernelMesh.kernel\_tendsto} proves pointwise convergence to a supplied
continuous kernel when the samples approach their points;
\texttt{KernelMesh.energy\_tendsto} passes the limit through the integral under
compactness, second countability and finite measure. The microscopic actuator
provides one such kernel from declared mode profiles and occupancies. These
results are in \nolinkurl{Phase2_KernelMesh.lean} and
\nolinkurl{Phase3_LocalActuator.lean}. The cell-pair integral identity and the
two convergence results are included in the extracted appendix.
\formal{cell-kernel-integral}
\formal{cell-kernel-limit}
\formal{cell-energy-limit}

\scope The vertex-supported no-go does not obstruct this cell-pair construction.
Its limiting kernel is specified by the hardware model, rather than inferred
from a scalar energy alone. The chain receives a concrete field through $E56$,
whose kernel conditions constrain the admissible input. The downstream
fixed-point proof uses only the scalar part of $E56$; the field conditions are
consumed separately by \texttt{em\_field\_exhibits\_phase\_transition}.
Identifying that field with cortex remains a physical modelling obligation.
